% ============================================================
%  CAPÍTULO 5 — ARQUITECTURA FÍSICA (Subdoc. 4.2)
%  Contenido: infraestructura, hardware, red, nube, on-premise
% ============================================================
%  FUENTES:
%  - Subdoc. 4 (Bases_Administrativas.md §1556-1565), parte física:
%      emplazamiento nube/on-premise justificado por componente, Art. 16° (§1559).
%      despliegue: ambientes, HA, DR, respaldos (§1562).
%      dimensionamiento y plan de capacidad (§1563).
%  - Híbrido obligatorio: AA Art. 16° y TT §4.1.
%  - Emplazamiento on-premise del caso: C §738 (RT-06.01: sala ténica CD Talca
%    25 m2; gabinetes de borde CD Concepción + 3 plataformas cross-docking).
%  - Operación desconectada: C §731 (RT-03.10: CD 24h sin enlace, dispositivos 14h).
%  - Ambientes / DR: AA Art. 24°, AA Art. 20°, TT RT-04.01, RT-07.07 (ver GG #3).
%  Ver FUENTES.md para el detalle completo.
\chapter{Arquitectura Física (Subdoc.\ 4.2)}

\nota{Este capítulo describe dónde y sobre qué infraestructura vive la solución: servidores, almacenamiento, red, dispositivos móviles, componentes IoT. Las Bases Técnicas Transversales (Art. 4.1) exigen despliegue híbrido obligatorio: componente en la nube + componente on-premise. Este capítulo debe demostrar que el despliegue híbrido es viable, seguro y está correctamente dimensionado.}

% -----------------------------------------------------------
\section{Estrategia de Despliegue Híbrido}
% -----------------------------------------------------------

\nota{Justificar por qué se usa despliegue híbrido (nube + on-premise) y no solo nube o solo on-premise. En el contexto de Puelche: (1) el centro de distribución en Concepción puede perder conectividad y debe seguir operando (on-premise); (2) la analítica y el dashboard de costo de servir escalan mejor en la nube; (3) la Ley 20.900 puede exigir que ciertos datos residan en territorio nacional. Incluir un diagrama de alto nivel que muestre qué componentes van en nube y cuáles on-premise.}

\pendiente

% -----------------------------------------------------------
\section{Componente On-Premise}
% -----------------------------------------------------------

\nota{Describir la infraestructura en el centro de distribución de Concepción. Este componente debe garantizar operación continua del centro de distribución (mínimo 24 horas continuas de recepción, preparación y despacho) sin enlace de red, según RT-03.10.}

\subsection{Servidor Local (Edge)}
\nota{Especificaciones del servidor on-premise: CPU, RAM, almacenamiento (SSD/NVMe), sistema operativo, virtualización (Proxmox, VMware, o bare metal). Justificar el dimensionamiento según la carga esperada (recepciones, preparación, despacho simultáneos).}

\pendiente

\subsection{Dispositivos de Terreno}
\nota{Dispositivos para conductores y preventistas: tablets o smartphones, escáneres de código de barras/QR, dispositivos IoT para temperatura (sensores Bluetooth o LoRa). Deben operar un turno completo de 14 horas sin cobertura (RT-03.10). Incluir especificaciones mínimas y justificación.}

\pendiente

\subsection{Red Local}
\nota{Infraestructura de red dentro del centro de distribución: switches, access points WiFi, segmentación de red (VLANs), cortafuegos locales. Considerar: cobertura WiFi en bodega, carga de dispositivos, y respaldo de conectividad (4G/5G).}

\pendiente

\subsection{Alimentación Eléctrica y Continuidad}
\nota{UPS (baterías de respaldo), generador eléctrico si aplica, plan de continuidad ante cortes de energía. El centro de distribución opera 24/7 y no puede detenerse por falta de energía.}

\pendiente

% -----------------------------------------------------------
\section{Componente en la Nube}
% -----------------------------------------------------------

\nota{Describir la infraestructura en la nube (AWS, Azure, o GCP). Justificar la selección del proveedor: costo, disponibilidad en Chile (regiones cercanas), cumplimiento normativo, ecosistema de servicios.}

\subsection{Proveedor de Nube}
\nota{Comparar al menos 2 proveedores (AWS vs Azure vs GCP) y justificar la elección. Considerar: región más cercana a Chile, costo estimado mensual, servicios relevantes (compute, storage, managed databases, IoT hub).}

\pendiente

\subsection{Servicios Utilizados}
\nota{Listar los servicios específicos de nube que se usarán: EC2/VM para compute, RDS/Cloud SQL para base de datos, S3/Blob Storage para archivos, IoT Hub para sensores, CDN para contenido estático. Para cada uno: tamaño estimado, costo mensual aproximado, y justificación.}

\pendiente

\subsection{Disaster Recovery}
\nota{Plan de recuperación ante desastres: backup automático (frecuencia, retención), réplica cross-region si aplica, RTO (Recovery Time Objective) y RPO (Recovery Point Objective) definidos. Debe ser consistente con los acuerdos de nivel de servicio del Art. 4.1.}

\pendiente

% -----------------------------------------------------------
\section{Conectividad entre Componentes}
% -----------------------------------------------------------

\nota{Diagrama de red que muestre cómo se conectan los componentes: centro de distribución $\leftrightarrow$ nube $\leftrightarrow$ dispositivos móviles. Protocolos de comunicación: MQTT para IoT, HTTPS para APIs, WebSockets para actualizaciones en tiempo real. Considerar: latencia tolerable, ancho de banda, y comportamiento ante desconexión.}

\pendiente

% -----------------------------------------------------------
\section{Dimensionamiento y Costos de Infraestructura}
% -----------------------------------------------------------

\nota{Tabla consolidada del dimensionamiento de infraestructura: componentes on-premise (qty × costo unitario), servicios de nube (qty × costo mensual), dispositivos móviles (qty × costo unitario), y costo total estimado. Incluir comparación con costo actual de la operación (si los datos del caso lo permiten).}

\begin{tabularx}{\textwidth}{|L{3cm}|C{1.5cm}|L{2.5cm}|L{2.5cm}|X|}
\hline
\rowcolor{lafroxClaro}
\textbf{Componente} & \textbf{Qty} & \textbf{Costo unitario} & \textbf{Costo total} & \textbf{Justificación} \\
\hline
Servidor on-premise & 1 & \pendiente & \pendiente & Edge computing para CD \\
\hline
Tablets/dispositivos & \pendiente & \pendiente & \pendiente & Conductores y preventistas \\
\hline
Sensores temperatura & \pendiente & \pendiente & \pendiente & Cadena de frío \\
\hline
Servidor nube (mensual) & \pendiente & \pendiente & \pendiente & Analítica, dashboard \\
\hline
Base de datos nube & \pendiente & \pendiente & \pendiente & Datos transaccionales \\
\hline
\end{tabularx}

% -----------------------------------------------------------
\section{Seguridad Física y Lógica}
% -----------------------------------------------------------

\nota{Medidas de seguridad en cada capa: acceso físico al servidor on-premise (candado, registro), acceso lógico (autenticación, autorización, cifrado en tránsito y en reposo), gestión de secretos (vault), parcheo, y auditoría. Alinear con ISO 27001 si se menciona en las credenciales del Cap. 1.}

\pendiente
